\documentclass[10pt,a4paper]{amsart}
\usepackage[margin=19mm]{geometry}
\usepackage{amsmath,amssymb,mathtools}
\usepackage[scr=rsfs,cal=euler]{mathalfa}
\usepackage{xfrac}
\usepackage[unicode,hidelinks,bookmarks=false]{hyperref}
\newcommand{\ie}{i.e.\ }
\newcommand{\cf}{cf.\ }
\newcommand{\resp}{respectively\ }
\newcommand{\Subsection}{\subsection}
\providecommand{\nobreakdash}{\nobreak\hbox{-}}
\setlength{\emergencystretch}{2em}
\multlinegap0pt
\usepackage{tikz}
\usetikzlibrary{cd}
\usepackage{xcolor}
\usepackage{amssymb}
\usepackage{mathtools}
\usepackage[shortlabels]{enumitem}
\usepackage{tikz}
\newtheorem{lemma}{Lemma}
\newtheorem{prop}{Proposition}
\newcommand{\F}{\mathbb{F}}
\newcommand\N{\mathbb{N}}
\newcommand\Sym{{\mathrm {Sym}}}
\newcommand{\Z}{\mathbb{Z}}
\newcommand\sF{{\mathcal{F}}}
\title{Behaviour of certain crystalline representations modulo $2$}
\author{Shalini Bhattacharya, Arathy Venugopal}
\date{Source: arXiv:2605.03424v1 (CC BY 4.0)}
\begin{document}
\maketitle

\noindent\emph{Source and licence.} Behaviour of certain crystalline representations modulo $2$. Version arXiv:2605.03424v1 (CC BY 4.0), distributed by arXiv under the Creative Commons Attribution 4.0 International licence, http://creativecommons.org/licenses/by/4.0/, as recorded in the arXiv licence metadata for this exact version and shown on the abstract page https://arxiv.org/abs/2605.03424v1. Full licence notice: \texttt{licenses/arxiv-2605.03424-CC-BY-4.0.txt}.\par\smallskip

\noindent\textit{Editorial scope.} Counted unit: the article's prop F2' (source0.tex lines 2120--2125). The article's complete original proof of it is delivered with it (source0.tex lines 2126--2320, one \texttt{proof} environment of 195 lines). No mathematics is written, repaired or re-derived: the statement and the proof are the article's own lines, in the article's own notation, and no retained line is substituted or re-flowed.  Retained uncounted as the source-local context the proof needs: lines 411--426; lines 540--576; lines 1555--1574; lines 1712--1717; lines 1960--1965; lines 1980--1991; lines 2026--2042.  Nothing was omitted from any retained span, so the package carries no omission record.  No retained line contains a citation, so the wrapper carries no bibliography and no bibliography entry.  The retained lines contain the article's own diagram code, which the wrapper typesets with the standard tikz packages; no image file is delivered and the package declares no asset dependency.  Editorial formatting only, in addition: an AMS \texttt{amsart} preamble that restates the article's own declarations and macros used by the retained lines, namely the environments \texttt{lemma}, \texttt{prop} on separate counters, together with the standard packages those lines need.  The retained environments print in this document as Lemma 1, Proposition 1 in order of appearance, whereas the article prints the same items with the numbering of its own full document.  The complete native source of the article as distributed in the arXiv e-print for this exact version is bundled unchanged as \texttt{sources/199789/source0.tex}.
\begin{lemma}\label{T}
 Let $P(X,Y) \in \Sym^{r}R^{2}$ where $R$ is a $\mathbb{Z}_{2}$-algebra. Let $\lambda \in I_{n},\,n\geq 0$. 
 Then the action of $T=T^+ + T^-$ on $\mathcal{I}(\Sym^{r}(R^{2}))$ can be explicitly stated as follows:
\begin{eqnarray*}
T^+\left[g_{n,\lambda}^0,\enspace P(X,Y)\right]& = &\left[g_{n+1,\lambda}^{0},\enspace P(X,2Y)\right] +  \left[g_{n+1, \lambda+2^n}^{0},\enspace  P(X,2Y-X)\right] \,\forall \,n\geq0,\\
T^-\left[g_{n,\lambda}^0,\enspace P(X,Y)\right] &= & \begin{cases} \left[g_{n-1,\lambda}^{0}, \enspace P(2X,Y)\right]& \text{ if } \lambda \in I_{n-1}\\
\left[g_{n-1,[\lambda]_{n-1}}^{0},\enspace P(2X, X+Y)\right]& \text{ if } \lambda \in I_{n}\setminus I_{n-1}
\end{cases}\forall\, n\geq1,\\
T^+\left[g_{n,\lambda}^1,\enspace P(X,Y)\right]& = &\left[g_{n+1,\lambda}^{1},\enspace P(2X,Y)\right] +  \left[g_{n+1, \lambda+2^n}^{1},\enspace P(2X-Y,Y)\right] \,\forall\,n\geq0,\\
T^-\left[g_{n,\lambda}^1,\enspace P(X,Y)\right] &= & \begin{cases} \left[g_{n-1,\lambda}^{1},\enspace P(X,2Y)\right]& \text{ if } \lambda \in I_{n-1}\\
\left[g_{n-1,[\lambda]_{n-1}}^{1},\enspace P(X+Y, 2Y)\right]& \text{ if } \lambda \in I_{n}\setminus I_{n-1}
\end{cases}\, \forall\, n\geq1,\\
T^-\left[g_{0,0}^0,\enspace P(X,Y)\right] &= & \left[g_{0,0}^{1},\enspace P(2X,Y)\right],\\
T^-\left[g_{0,0}^1,\enspace P(X,Y)\right] &= & \left[g_{0,0}^{0},\enspace  P(X,2Y)\right].
\end{eqnarray*}
\end{lemma}

\begin{lemma}\label{T+toT-}
    Let $n \in \N \cup \{0\}$ and $\lambda \in I_{n}$, $w=\left(\begin{smallmatrix}
        0 & 1\\
        1 & 0
    \end{smallmatrix}\right)$ and let $P(X,Y)\in \Sym^rR^2$ where $R$ is a $\Z_{2}$-algebra. Then we have:  \begin{eqnarray}T^+\left[g_{n,\lambda}^{1}, \enspace w \cdot P(X,Y)\right]&=&\left(\begin{smallmatrix}
    0 & 1\\
    1 & \lambda
\end{smallmatrix}\right)\cdot T^+\left[g_{n+1,\,\lambda}^{0},\enspace P(X,Y)\right]
   %\\ &=& \left[g_{n+1,\lambda}^{1}, \textcolor{red}{w.S(X,2Y)}\right] + \left[g_{n+1,\lambda+2^n}^{1}, w.S(X,2Y-X)\right]
    %where $\left[g_{n+2,\lambda}^{0},S(X,Y)\right] \to \left[g_{n+1,\lambda}^{1}, w.S(X,Y)\right]$ and $\left[g_{n+2, 2^{n+1}+\lambda}^{0},S(X,Y)\right] \to \left[g_{n+1,2^{n}+\lambda}^{1}, w.S(X,Y)\right]$ under
    %the action of $\begin{pmatrix}
    %    0 & 1\\
     %   1 & \lambda
   % \end{pmatrix}.$ 
   \label{T+toT-a}\\
   % \item $ T^+\left[g_{n,\lambda}^{1}, w.S(X,Y)\right]
   &=&\left(\begin{smallmatrix}
        0 & 1\\
        1 & \lambda -2^n
\end{smallmatrix}\right)\cdot T^+\left[g_{n+1,\, \lambda+2^n}^{0},\enspace P(X,Y)\right].\label{T+toT-b}\\
    %,$ where $\left[g_{n+2,2^n+\lambda}^{0},S(X,Y)\right] \to \left[g_{n+1,\lambda}^{1}, w.S(X,Y)\right]$ and $\left[g_{n+2, 2^{n+1}+2^n+\lambda}^{0},S(X,Y)\right] \to \left[g_{n+1,2^{n}+\lambda}^{1}, w.P(X,Y)\right]$ under
    %the action of $\begin{pmatrix}
    %    0 & 1\\
    %    1 & \lambda-2^n
%\end{pmatrix}.$ 
T^-\left[g_{n,\lambda}^1, \enspace w \cdot P(X,Y)\right] &=& \left(\begin{smallmatrix}
       0&1\\
       1&\lambda
   \end{smallmatrix}\right)\cdot T^-\left[g_{n+1,\,\lambda}^{0}, \enspace P(X,Y)\right]
   \label{T+toT-c}\\
   %T^-\left[g_{n, \lambda}^{1}, w.P(X,Y)\right]
   &=& \left(\begin{smallmatrix}
       0 & 1\\
       1 & \lambda-2^n
   \end{smallmatrix}\right)\cdot T^-\left[g_{n+1,\,\lambda+2^n}^{0},\enspace P(X,Y)\right]. \label{T+toT-d}
  \end{eqnarray} 
    \end{lemma}

\begin{lemma}\label{Heven}
Let $r\geq 3$ and $t=v(r-2)$. Then modulo $2^{t+2}$, we have, 
\begin{enumerate}
	\item $H(X,2Y) \equiv 
    \begin{cases}
    2rX^{r-1}Y + 2rX^{r-2}Y^2 &\text{ if }\, r \text{ is even}\\
    2rX^{r-1}Y &\text{ if }\, r \text{ is odd}    
    \end{cases}
    $ \label{Heven1}
	\item $H(X,2Y-X) \equiv \begin{cases}
-2X^{r}+2rX^{r-1}Y-2rX^{r-2}Y^2 & \text{ if }\, r \text{ is even,}\\
 -2rX^{r-1}Y & \text{ if }\, r \text{ is odd.}
    \end{cases}$    
    \label{Heven2}
	\item $H(2X,Y) \equiv \begin{cases}
2rXY^{r-1} + 2rX^2Y^{r-2} \text{ if }\, r \text{ is even,}\\
2rXY^{r-1} \text{ if }\, r \text{ is odd.}
	\end{cases}$ \label{Heven3}
\end{enumerate}
\end{lemma}

\begin{equation}\label{diag4}
\begin{tikzcd}
 		0 \arrow[r] & \mathcal{I}\left(V_{1}\right) \arrow[d, ""]\arrow[r] & \mathcal{I}\left(\dfrac{V_{r}^{(1)}}{V_{r}^{(2)}}\right)  \arrow[r] \arrow[d, "P"] & \mathcal{I}(V_{0}) \arrow [r] \arrow[d, ""]& 0\\
 		0 \arrow[r] &\mathcal{F'}_{2}  \arrow[r] & \bar{\Theta}_{0} \arrow[r] & \mathcal{F'}_{3}\arrow[r] & 0
 \end{tikzcd}
\end{equation}

\begin{lemma}\label{sl1par}  \begin{enumerate}
      \item The quantity $\frac{a_{2}^2-2r}{a_{2}^2\beta_{0}}$ is always integral and 
   it is a unit congruent to $\frac{2}{a_{2}}\mod\wp$ when $\tau\leq t-1$. when $\tau>t-1$, it vanishes if $r$ is even and otherwise a unit congruent to $\frac{2^2}{a_{2}^2}\mod\wp$. \label{sl1par1} 
   \item The quantity $\frac{a_{2}^2-2^2}{2a_{2}\beta_{0}}$ is always integral and is a congruent to $\left(1+\frac{(r-2)}{a_{2}\beta_{0}}\right)$ when $\tau\leq t-1$. When $\tau>t-1$, then it vanishes mod $\wp$ if $r$ is odd and when $r$ is even, it is a unit congruent to $\frac{2}{a_{2}} \mod\wp$.\label{sl1par2} 
   \end{enumerate} 
\end{lemma}

\begin{lemma}\label{sl1cmb2}
Let $r\geq 4$ and $t=v(r-2)$. Then  we have  
\begin{enumerate}
    \item 
$v\left({r-1 \choose j}\right)+j \geq t+1$ for all $j\geq2$, and \label{sl1cmb2a}
\item $v\left({r-2 \choose j}\right)+j \geq t+1$ for all $j\geq1$.\label{sl1cmb2b}
\end{enumerate}
%In particular, 
%$$ \nu\left({r-1 \choose j}\right)+j \begin{cases}
%%	t+1 \,\,\text{for} \, j=2,3.
%\end{cases}$$
\end{lemma}

\begin{lemma}\label{K}
For $r\geq 6$, let $K(X,Y):= \displaystyle\sum_{j=3}^{r-3}\left[{r-1 \choose j}+{r-2\choose j} \right]X^{r-j}Y^j$. %\text{ where } r\geq6 \text{ even}$. 
Let $t:=v(r-2)$. Then we have,
\begin{enumerate}
    \item $K(X,2Y) \equiv \begin{cases}
        0\mod 2^{t+1} & \text{ if }r\text{ even},\\
        0\mod 2^2 & \text{ if }r\text{ odd }.
        \end{cases} $
    \item $K(X,2Y-X) \equiv\begin{cases}
    -rX^{r}-2^2X^{r-1}Y\mod 2^{t+1} & \text{if $r$ even,}\label{Kb}\\
     -2{r-1 \choose 2}X^{r} \mod 2^{2}& \text{if $r$ odd.}    \end{cases}$
    \item $K(2X,Y) \equiv \begin{cases}
        0\mod 2^{t+1} & \text{ if }r\text{ even},\\
        0\mod 2^2 & \text{ if }r\text{ odd }.
        \end{cases} $
    \end{enumerate}
\end{lemma}

\begin{prop}\label{F2'}
    Let $r\geq6$ be even, $v(a_{2})=1$ and let $\tau$ and $t$ be as above.
    \begin{enumerate}
        \item If $\tau\leq t-1$, then $\mathcal{F'}_{2}$ is a quotient of $\dfrac{\mathcal{I}(V_{1})}{(T^{2}-c_{0}T+1)}$ where $c_{0} = \overline{\dfrac{a_{2}}{2} + \dfrac{r-2}{2\alpha}}$ .
        \item If $\tau>t-1$, then $\mathcal{F'}_{2}$ is a quotient of $\dfrac{\mathcal{I}(V_{1})}{(T)}$.                             \end{enumerate}
\end{prop}

\begin{proof}
Consider function $f=\underset{n\in\Z}\sum f_n$, where $f_n$'s are defined as follows:
\begin{eqnarray}
f_{0}&=& \left[1, \frac{X^{r}+X^{r-1}Y+X^{2}Y^{r-2}}{a_{2}\beta_{0}} + \frac{2r}{a_{2}^{3}\beta_{0}}K(X,Y)\right]\nonumber\\
%f_{2} &= &\left[g_{2,0}^{0}, \dfrac{Y^{r} - X^{r-2}Y^{2}}{a_{2}\beta_{0}}\right]  - \left[g_{2,2}^{0}, \dfrac{Y^{r} - X^{r-2}Y^{2}}{a_{2}\beta_{0}}\right] - \left[g_{2,1}^{0}, \dfrac{Y^{r} - X^{r-2}Y^{2}}{a_{2}\beta_{0}}\right] \nonumber\\
%& + & \left[g_{2,3}^{0}, \dfrac{Y^{r} - X^{r-2}Y^{2}}{a_{2}\beta_{0}}\right]\textcolor{red}{\text{can be written along with fn}}\nonumber\\
f_{1} &=& \left[g_{1,0}^{0}, \dfrac{-H(X,Y)}{a_{2}^{2}\beta_{0}}+ \dfrac{X^{r-1}Y}{2\beta_{0}}+ \dfrac{c_{2}}{a_{2}}X^{r-2}Y^{2}\right]\label{F2'eqeven}\\
& & - \left[g_{1,1}^{0},  \dfrac{-H(X,Y)}{a_{2}^{2}\beta_{0}}+ \dfrac{X^{r-1}Y}{2\beta_{0}}+  \dfrac{c_{2}}{a_{2}}X^{r-2}Y^{2}                \right]  \nonumber\\
f_{n} &=& \left[g_{n,0}^{0}, \dfrac{a_{2}^{n-3}}{\beta_{0}}(Y^{r}-X^{r-2}Y^{2})\right] -  \left[g_{n,2}^{0}, \dfrac{a_{2}^{n-3}}{\beta_{0}}(Y^{r}-X^{r-2}Y^{2})\right] \nonumber\\ 
& & - \left[g_{n,1}^{0}, \dfrac{a_{2}^{n-3}}{\beta_{0}}(Y^{r}-X^{r-2}Y^{2})\right] +  \left[g_{n,3}^{0}, \dfrac{a_{2}^{n-3}}{\beta_{0}}(Y^{r}-X^{r-2}Y^{2})\right]\text{ for } 2\leq n \leq t+2 \nonumber\\
f_{-n} &= &\left[g_{n-1,0}^{1}, -\dfrac{a_{2}^{n-1}}{\beta_{0}}(X^{r}-X^{2}Y^{r-2})\right] \text{ for }\,1 \leq n\leq t,\nonumber
%f_{-(n+1)} &= &\left[g_{n,0}^{1}, -\dfrac{2^{2n+2}}{a_{2}^{n+2}\beta_{0}}(X^{r}-X^{2}Y^{r-2})\right] \text{for all}\, n\geq0. \nonumber
\end{eqnarray}
where $K(X,Y)$ is as defined in Lemma \ref{K} and $c_{2}$ is as in Lemma \ref{sl1par}.
We now compute  $(T-a_{2})f$ radius by radius.

\noindent\underline{\textbf{Radius 0}}\\
 Applying Lemma \ref{T}, we get %the definition of $T$ in L.H.S of the tree, we can compute $T^{-}f_{-1}$ modulo $\wp$ as,
\begin{equation}\label{F2'T-f-1}
    T^{-}f_{-1} \equiv 
        \left[1, \dfrac{-X^r}{\beta_{0}}\right]\mod\wp,
\end{equation}
since $v(\beta_0)\leq t-1<r-2$.
Further we have ,
\begin{equation}\label{F2'a2f0}
 -a_{2}f_{0} = \left[1, \dfrac{-\left(X^{r}+X^{r-1}Y+X^{2}Y^{r-2}\right)}{\beta_{0}} -\dfrac{2r}{a_{2}^2\beta_{0}}K(X,Y) \right].
 \end{equation}
We now calculate $T^{-}f_{1}$ using Lemma \ref{Heven} and the condition that $r$ is even:
\begin{eqnarray}\label{F2'T-f1}      T^{-}f_1 &\equiv& \left[1, \dfrac{-H(2X,Y)}{a_{2}^2\beta_{0}}+\dfrac{H(2X,X+Y)}{a_{2}^2\beta_{0}}\right] \mod \wp \nonumber\\
    &\equiv&  \left[1, -\dfrac{(2rXY^{r-1}+2r    X^{2}Y^{r-2})}{a_{2}^2\beta_{0}}+\dfrac{2rX(X+Y)^{r-1}+2rX^{2}(X+Y)^{r-2}} {a_{2}^2\beta_{0}}\right]\mod\wp \nonumber\\
    &=&\left[1, \frac{2r}{a_2^2\beta_0}\left(2X^r+(2r-3)X^{r-1}Y+(r-2)^2X^{r-2}Y^2+K(X,Y)+(r-1)X^2Y^{r-2}\right)\right] \nonumber\\
    &\equiv&\left[1, \dfrac{2^2rX^r}{a_{2}^2\beta_{0}} + \dfrac{2r(X^{r-1}Y+X^{2}Y^{r-2})}{a_{2}^2\beta_{0}}+\dfrac{2r}{a_{2}^2\beta_{0}}K(X,Y)\right]\mod\wp.
\end{eqnarray}
The last congruence follows reducing the previous expression modulo $\wp$, using the fact that for $r$ even, the valuation $v\left(\frac{2r(r-2)}{a_2^2\beta_0}\right)= t-1-v(\beta_0)+v(r)\geq v(r)>0$.
%We note that the above congruence is obtained by isolating $K(X,Y)$ from $H(2X,X+Y)$, combining the remaining like terms in $T^-f_{1}$ and evaluating the coefficients modulo $t+1$.



 
 %\equiv\begin{cases}

 
%\left[1, \dfrac{-\left(X^{r}+X^{r-1}Y+X^{2}Y^{r-2}\right)}{\alpha} -\dfrac{2r}{a_{2}^2\alpha}\displaystyle\sum_{j=3}^{r-3}\left[{r-1 \choose j}+{r-2\choose j} \right]X^{r-j}Y^j   \right]& \text{ if } \tau\leq t-1\\
%0               & \text{ if } \tau>t-1
% \end{cases}&
%\end{eqnarray}
 On adding \eqref{F2'T-f-1},\eqref{F2'a2f0} and \eqref{F2'T-f1}, we get that,
 \begin{eqnarray}
T^{-}f_{-1}-a_{2}f_{0}+T^{-}f_{1} &\equiv& \left[1, -\dfrac{(a_{2}^2-2r)}{a_{2}^2\beta_{0}}(X^{r-1}Y+X^2Y^{r-2}+2X^r)\right] \nonumber\\
&\equiv& \left[1, -c_{1}(X^{r-1}Y-X^2Y^{r-2})\right]\mod\wp \nonumber
\end{eqnarray}
 since $c_{1}$ is integral by Lemma \ref{sl1par}.  
 %&\equiv& \left[1,\dfrac{(-a_{2}^2+2r)X^{r-3}\theta}{a_{2}^2\beta_{0}} -\dfrac{(a_{2}^2-2^2)X^r}{a_{2}^2\beta_{0}}     \right]\mod\wp.  
By using the method of alternating sum, we note that $X^{r-1}Y-X^{2}Y^{r-2} \equiv \theta X^{r-3} + \theta \displaystyle\sum_{j=1}^{r-4}X^{r-3-j}Y^{j} \equiv \theta X^{r-3} \mod V_{r}^{(2)}$ since $r$ is even. 
\begin{equation}\label{sl1F2'rad0}
T^{-}f_{-1}-a_{2}f_{0}+T^{-}f_{1} \equiv \begin{cases} 
\left[   1, -\dfrac{2}{a_{2}}\theta X^{r-3}\right] \mod(\wp, V_{r}^{(2)}+X^{r}) & \text{ if }\tau\leq t-1,\\
0 \mod(\wp, V_{r}^{(2)}+X^{r}) & \text{ if } \tau>t-1. 
\end{cases}
\end{equation}
\noindent\underline{\textbf{Radius 1}}\\
%For computing $T^{+}f_{0}$, we shall use 
Using Lemma \ref{sl1cmb2} and Lemma \ref{K} along with the definition of $T^+$, we compute
\begin{eqnarray}\label{F2'T+f0}
    T^{+}f_{0} &\equiv& \left[g_{1,0}^{0}, \dfrac{X^{r}+ 2X^{r-1}Y}{a_{2}\beta_{0}}\right] + \left[g_{1,1}^{0}, \dfrac{X^{r}+2X^{r-1}Y}{a_{2}\beta_{0}}+\dfrac{2r}{a_{2}^{3}\beta_{0}}K(X,2Y-X) \right] \nonumber\\
    &\equiv& \left[g_{1,0}^{0},\dfrac{X^{r}+ 2X^{r-1}Y}{a_{2}\beta_{0}}\right]+ \left[g_{1,1}^{0}, \dfrac{X^{r}+ 2X^{r-1}Y}{a_{2}\beta_{0}} + \dfrac{-2r^2X^{r}-2^3rX^{r-1}Y}{a_{2}^3\beta_{0}}   \right] \nonumber\\
    &\equiv& \left[g_{1,0}^{0},\dfrac{X^{r}+ 2X^{r-1}Y}{a_{2}\beta_{0}}\right]+ \left[g_{1,1}^{0}, \dfrac{(a_{2}^2-2r^2)X^{r}}{a_{2}^3\beta_{0}} + \dfrac{(2a_{2}^2-2^3r)X^{r-1}Y}{a_{2}^3\beta_{0}}   \right] \nonumber\\
   &\equiv& \left[g_{1,0}^{0},\dfrac{X^{r}+ 2X^{r-1}Y}{a_{2}\beta_{0}}\right]+ \left[g_{1,1}^{0}, \dfrac{(a_{2}^2-2r^2)X^{r}}{a_{2}^3\beta_{0}} - \dfrac{2X^{r-1}Y}{a_{2}\beta_{0}}   \right]   \mod \wp.
    \end{eqnarray}
The last congruence follows from Lemma \ref{sl1par} part \eqref{sl1par1}. 
%We shall now simplify the coefficient of $X^{r-1}Y$ in the part of $T^{+}f_{0}$ supported on $g_{1,1}^{0}$ in the above congruence.
%\begin{equation*}
 %   \dfrac{(2a_{2}^2-2^3r)}{a_{2}^3\beta_{0}} \equiv \dfrac{-2a_{2}^{2}+2^2a_{2}^2-2^3r}{a_{2}^3\beta_{0}}\equiv \dfrac{-2a_{2}^2+2^2(a_{2}^2-2r)}
 %   {a_{2}^3\beta_{0}} 
 %   \equiv \dfrac{-2}{a_{2}\beta_{0}} \mod \wp
  %  \end{equation*}
% Thus we have \begin{equation}
%T^{+}f_{0}\equiv \left[g_{1,0}^{0},\dfrac{X^{r}+ 2X^{r-1}Y}{a_{2}\beta_{0}}\right]+ \left[g_{1,1}^{0}, \dfrac{(a_{2}^2-2r^2)X^{r}}{a_{2}^3\beta_{0}} -\dfrac{2X^{r-1}Y}{a_{2}\beta_{0}}\right].
%\end{equation}
Next we have
\begin{eqnarray}\label{F2'a2f1}
    -a_{2}f_{1} &=& \left[g_{1,0}^{0}, \dfrac{H(X,Y)}{a_{2}\beta_{0}}-\dfrac{a_{2}X^{r-1}Y}{2\beta_{0}} -c_{2}X^{r-2}Y^{2}\right] \nonumber\\
    &&+ \left[g_{1,1}^{0},-\dfrac{H(X,Y)}{a_{2}\beta_{0}}+\dfrac{a_{2}X^{r-1}Y}{2\beta_{0}} +c_{2}X^{r-2}Y^{2}\right]. 
\end{eqnarray}
%We now compute $T^{-}f_{2}$ 
Now using the definition of $T^-$, we compute
\begin{equation}\label{F2'T-f2}
    T^{-}f_{2} = \left[g_{1,0}^{0}, \dfrac{-X^{r}-H(X,Y)}{a_{2}\beta_{0}}\right] + \left[g_{1,1}^{1}, \dfrac{X^{r}+H(X,Y)}{a_{2}\beta_{0}}\right].
\end{equation}

On adding up \eqref{F2'T+f0},\eqref{F2'a2f1} and \eqref{F2'T-f2} we get modulo $\wp$,
\begin{equation*}
  T^{+}f_{0}-a_{2}f_{1}+T^{-}f_{2}\equiv \left[g_{1,0}^{0}, -c_{2}\theta X^{r-3}\right] + \left[g_{1,1}^{0}, \dfrac{2(a_{2}^2-r^2)X^{r}}{a_{2}^3\beta_{0}}+ c_{2}\theta X^{r-3} \right].
    \end{equation*}
     The coefficient of $X^{r}$ above is congruent to $\frac{2^2a_{2}\alpha-2r(r-2)}{a_{2}^3\beta_{0}} = \frac{2^2\alpha}{a_{2}^2\beta_{0}}-\frac{2r(r-2)}{a_{2}^3\beta_{0}} \mod \wp$. Since $v(\beta_{0})= \min(\tau,t-1)$, this coefficient is integral. Thus going modulo $(\wp,V_{r}^{(2)}+X_{r})$ and further applying  Lemma \ref{sl1par} \eqref{sl1par2} for $r$  even, we have 
    \begin{eqnarray}\label{sl1F2'rad1}
   T^{+}f_{0}-a_{2}f_{1}+T^{-}f_{2}
    &\equiv &\begin{cases} 
    \left[g_{1,0}^{0}, -\left(1+\frac{r-2}{a_{2}\alpha} \right)\theta X^{r-3}\right] + \left[g_{1,1}^{0}, \left(1+\frac{r-2}{a_{2}\alpha}\right)\theta X^{r-3} \right] & \text{ if } \tau\leq t-1,\\
     \left[g_{1,0}^{0}, \frac{-2\theta X^{r-3}}{a_{2}}\right] + \left[g_{1,1}^{0}, \frac{2\theta X^{r-3}}{a_{2}}\right]  &\text{ if } \tau>t-1.
     \end{cases}\nonumber\\
     &&
     \end{eqnarray}
 %The above congruence follows from Lemma \ref{sl1par} \eqref{sl1par2} and the fact that $r$ is even. 
\noindent\underline{\textbf{Radius $n\geq 2$}}\\
We now compute $T^{+}f_{n-1}$ for $n\geq 3.$ In  this computation, the second branch of $T^{+}$ with the action $(X,Y)\to(X,2Y-X)$ vanishes  by \ref{sl1cmb2}\eqref{sl1cmb2b}, and we have
\begin{eqnarray}\label{F2'T+fn-1}
  T^{+}f_{n-1} &\equiv& \left[g_{n,0}^{0}, \dfrac{-2^2a_{2}^{n-4}X^{r-2}Y^2}{\beta_{0}}\right] +\left[g_{n,2}^{0}, \dfrac{2^2a_{2}^{n-4}X^{r-2}Y^2}{\beta_{0}}\right]\nonumber\\
  &&+\left[g_{n,1}^{0}, \dfrac{2^2a_{2}^{n-4}X^{r-2}Y^2}{\beta_{0}}\right] +\left[g_{n,3}^{0}, \dfrac{-2^2a_{2}^{n-4}X^{r-2}Y^2}{\beta_{0}}\right] \mod\wp.
\end{eqnarray}
%We note that in the \textcolor{red}{computation of $T^+$ above}, the second branch of $T^{+}$ with the action $(X,Y)\to(X,2Y-X)$ vanishes and it can be proved by applying Lemma \ref{sl1cmb2} part \eqref{sl1cmb2b}.
%
%\left[g_{n,2^{n-1}}, \dfrac{a_{2}^{n-4}(2Y-X)^{2}(2^j\sum_{j=1}^{r-2}X^{r-2-j}Y^{j})}{\beta_{0}}\right]

However for n=2, the calculation is more involved and we have
\begin{eqnarray*}
    T^{+}f_{1} &=& \left[g_{2,0}^{0}, \dfrac{-H(X,2Y)}{a_{2}^2\beta_{0}}+\dfrac{X^{r-1}Y}{\beta_{0}} +\frac{2^2c_{2}}{a_{2}}X^{r-2}Y^2\right]\\ 
    &&+ \left[g_{2,2}^{0}, \dfrac{-H(X,2Y-X)}{a_{2}^2\beta_{0}}+\dfrac{X^{r-1}(2Y-X)}{2\beta_{0}} +\dfrac{c_{2}}{a_{2}}X^{r-2}(2Y-X)^2\right] \\
    &&-\left[g_{2,1}^{0}, \dfrac{-H(X,2Y)}{a_{2}^2\beta_{0}}+\dfrac{X^{r-1}Y}{\beta_{0}}+\frac{2^2c_{2}}{a_2}X^{r-2}Y^2\right] \\ 
    &&-\left[g_{2,3}^{0}, \dfrac{-H(X,2Y-X)}{a_{2}^2\beta_{0}}+\dfrac{X^{r-1}(2Y-X)}{2\beta_{0}}+\dfrac{c_{2}}{a_{2}}X^{r-2}(2Y-X)^2\right]. 
\end{eqnarray*}
By using Lemma \ref{Heven} for $r$ even and part \eqref{sl1par2} of Lemma \ref{sl1par}, we get modulo $\wp$,
\begin{eqnarray}\label{F2'T+f1}
T^{+}f_{1}&\equiv &\left[ g_{2,0}^{0}, \left( \dfrac{a_{2}^2-2r}{a_{2}^2\beta_{0}}\right)X^{r-1}Y-\dfrac{2rX^{r-2}Y^2}{a_{2}^2\beta_{0}}  \right] + \left[ g_{2,2}^{0}, \left( \dfrac{a_{2}^2-2r}{a_{2}^2\beta_{0}}\right)X^{r-1}Y+\dfrac{2rX^{r-2}Y^2}{a_{2}^2\beta_{0}}  \right] \nonumber\\
 && - \left[ g_{2,1}^{0},  \left( \frac{a_{2}^2-2r}{a_{2}^2\beta_{0}}\right)X^{r-1}Y-\frac{2rX^{r-2}Y^2}{a_{2}^2\beta_{0}}\right] - \left[ g_{2,3}^{0},  \left( \frac{a_{2}^2-2r}{a_{2}^2\beta_{0}}\right)X^{r-1}Y+\frac{2rX^{r-2}Y^2}{a_{2}^2\beta_{0}}  \right] \nonumber\\
 && \equiv \left[ g_{2,0}^{0},  c_{1}X^{r-1}Y-\dfrac{2rX^{r-2}Y^2}{a_{2}^2\beta_{0}}  \right] + \left[ g_{2,2}^{0},  c_{1}X^{r-1}Y+\dfrac{2rX^{r-2}Y^2}{a_{2}^2\beta_{0}}  \right] \nonumber\\
 && - \left[ g_{2,1}^{0},   c_{1}X^{r-1}Y-\frac{2rX^{r-2}Y^2}{a_{2}^2\beta_{0}}\right] - \left[ g_{2,3}^{0}, c_{1}X^{r-1}Y+\frac{2rX^{r-2}Y^2}{a_{2}^2\beta_{0}}  \right].
 \end{eqnarray}
We note that in the second branch of $T^{+}$, the coefficients of the $X^{r}$ terms add up to zero.


For $n\geq 2$, we have,
\begin{eqnarray}\label{F2'a2fn}
-a_{2}f_{n}& = &\left[g_{n,0}^{0}, \dfrac{-a_{2}^{n-2}(Y^{r}-X^{r-2}Y^{2} )}{\beta_{0}}\right] +  \left[g_{n,2}^{0}, \dfrac{a_{2}^{n-2}(Y^{r}-X^{r-2}Y^{2})}{\beta_{0}}\right] \nonumber\\
&&+ \left[g_{n,1}^{0},\dfrac{a_{2}^{n-2}(Y^{r}-X^{r-2}Y^{2})}{\beta_{0}}\right] + \left[g_{n,3}^{0},\dfrac{-a_{2}^{n-2}(Y^{r}-X^{r-2}Y^{2})}{\beta_{0}}\right].
\end{eqnarray}%$$\overline{(T-a_{2})(f)} \mod V_{r}^{**} = \left(T^{2} - \left(\overline{\dfrac{2^{2}-a_{2}^{2}}{2a_{2}\alpha}}\right)T + 1 \right) [1,X]= \left(T^{2} - \left(\overline{\dfrac{a_{2}}{2}}\right)T + 1 \right)[1,X]$$ 
Using the definition of $T^-$ and that $r-2>t-1$ in this context, we get that for all $n\geq2$,
\begin{equation}\label{F2'T-fn+1}
T^{-}f_{n+1} \equiv \left[g_{n,0}^{0}, \dfrac{a_{2}^{n-2}Y^{r}}{\beta_{0}}\right] -\left[g_{n,2}^{0}, \dfrac{a_{2}^{n-2}Y^{r}}{\beta_{0}}\right]-\left[g_{n,1}^{0}, \dfrac{a_{2}^{n-2}Y^{r}}{\beta_{0}}\right]+\left[g_{n,3}^{0}, \dfrac{a_{2}^{n-2}Y^{r}}{\beta_{0}}\right] \mod\wp.
\end{equation}
 Thus using \eqref{F2'a2fn}, \eqref{F2'T-fn+1}, \eqref{F2'T+fn-1} and \eqref{F2'T+f1}, we get that for $n\geq 3$, \begin{eqnarray}\label{sl1F2'radn}
& &T^{+}f_{n-1}-a_{2}f_{n}+T^{-}f_{n+1}  \nonumber \\
 &\equiv&
\left[g_{n,0}^{0},\dfrac {a_{2}^{n-4}(a_{2}^2-2^2)}{\beta_{0}}X^{r-2}Y^2 \right]
-\left[g_{n,2}^{0},\dfrac {a_{2}^{n-4}(a_{2}^2-2^2)}{\beta_{0}}X^{r-2}Y^2\right]  \nonumber\\
 & &  -\left[g_{n,1}^{0},\dfrac {a_{2}^{n-4}(a_{2}^2-2^2)}{\beta_{0}}X^{r-2}Y^2\right]+\left[g_{n,3}^{0},\dfrac {a_{2}^{n-4}(a_{2}^2-2^2)}{\beta_{0}}X^{r-2}Y^2\right]  \nonumber\\
 & \equiv& 2a_{2}^{n-3}c_{2}\left(\left[g_{n,0}^{0}, X^{r-2}Y^2 \right]-\left[g_{n,2}^{0}, X^{r-2}Y^2\right] -\left[g_{n,1}^{0}, X^{r-2}Y^2\right]+\left[g_{n,3}^{0}, X^{r-2}Y^2\right] \right)  \nonumber\\
& \equiv& 0 \mod \wp 
\end{eqnarray}
and for $n=2$,
\begin{eqnarray*}
 T^{+}f_{1}-a_{2}f_{2}+T^{-}f_{3}\equiv c_{1}\left(\left[g_{2,0}^0, \theta X^{r-3}\right] + \left[g_{2,2}^0, \theta X^{r-3}\right] -\left[g_{2,1}^0, \theta X^{r-3} \right]-\left[g_{2,2}^0, \theta X^{r-3} \right]\right).\enspace
\end{eqnarray*}
Thus going$\mod(\wp, X_{r}+V_{r}^{(2)})$, we have
\begin{equation}\label{sl1F2'rad2}   
    T^{+}f_{1}-a_{2}f_{2}+T^{-}f_{3}
\equiv\begin{cases} 
\dfrac{2}{a_{2}}\left([g_{2,0}^0, \theta X^{r-3}] + [g_{2,2}^0, \theta X^{r-3}]-[g_{2,1}^0, \theta X^{r-3}]-[g_{2,2}^0,X^{r-3}\theta]\right)&\text{ if } \tau\leq t-1\\
0 &\text{ if } \tau>t-1
\end{cases}
\end{equation}
by Lemma \ref{sl1par} \eqref{sl1par1} together with the fact that $r$ is even.

\noindent\underline{\textbf{Radius $-n$ for $n\geq 1$}}\\
We have
\begin{equation}\label{F2'a2f-n}
-a_{2}f_{-n}=\left[g_{n-1,0}^{1}, \dfrac{a_{2}^{n}(X^{r}-X^2Y^{r-2})}{\beta_{0}}\right]
\end{equation}
Using Lemma \ref{T+toT-} equation \eqref{T+toT-c} we have
\begin{equation}\label{F2'T-f-n-1}
    T^{-}f_{-(n+1)} \equiv \left[g_{n-1,0}^{1}, -\dfrac{a_{2}^{n}X^{r}}{\beta_{0}}\right]\mod\wp.
\end{equation}
We now compute $T^{+}f_{-(n-1)}$ for $n\geq2$ and $T^{-}f_{0}$. We have modulo $\wp$,
\begin{eqnarray}\label{F2'T+f-n+1}
    T^{+}f_{-(n-1)}&\equiv& \left[g_{n-1,0}^{1}, \dfrac{-a_{2}^{n-2}(2^rX^{r}-2^2X^2Y^{r-2})}{a_{2}^n\beta_{0}}\right] \equiv\left[g_{n-1,0}^{1}, \dfrac{2^2a_{2}^{n-2}X^2Y^{r-2}}{\beta_{0}}\right]\label{F2'T+f-n+1a} \\
    T^{-}f_{0}&\equiv& \left[g_{0,0}^{1}, \dfrac{2^2X^2Y^{r-2}}{a_{2}\beta_{0}}\right] \label{F2'T-f0}
\end{eqnarray}
We note that the second branch of $T^{+}$ in \eqref{F2'T+f-n+1a} vanishes modulo $\wp$ applying Lemma \ref{sl1cmb2} \eqref{sl1cmb2b}. Using \eqref{F2'a2f-n}, \eqref{F2'T-f-n-1}, \eqref{F2'T+f-n+1} and applying Lemma \ref{sl1par}\eqref{sl1par2}, we get that for $n\geq 2$,
\begin{eqnarray*}
-a_{2}f_{-n}+T^{-}f_{-(n+1)}+T^{+}f_{-(n-1)} \equiv \left[g_{n-1,0}^{1}, -2a_{2}^{n-1}c_{2}X^2Y^{r-2} \right] 
\equiv 0 \mod\wp
\end{eqnarray*}
and for $n=1$ further using \eqref{F2'T-f0},\begin{eqnarray*} 
-a_{2}f_{-1}+T^{-}f_{-2}+T^{-}f_{0} \equiv \left[g_{0,0}^{1}, -2c_2X^2Y^{r-2}\right]
\equiv 0 \mod\wp
\end{eqnarray*}
Thus radii $-n, n\geq1$ vanishes irrespective of the relation between $\tau$ and $t$. 

 We recall that in diagram \eqref{diag4}, the $\bar{\F}_{2}[G]$-module $\sF'_{2}$ is a quotient of $\mathcal{I}(V_{1})$ where $V_{1}$ is isomorphic to the sub-module $\dfrac{\theta X_{r-3}}{\theta X_{r-3}^{(1)}}$ of $\dfrac{V_{r}^{(1)}}{V_{r}^{(2)}}$. From the above computations we get that $(T-a_{2})f$ is integral. Its reduction modulo $(\wp, X_{r}+V_{r}^{(2)})$ is supported on radius $0,1$ and $2$ and in fact only on radius $1$ when $\tau>t-1$. Further we observe from equations    \eqref{sl1F2'rad0}, \eqref{sl1F2'rad1} and \eqref{sl1F2'rad2} that this reduction lands in the submodule $\mathcal{I}(V_{1})$ of diagram \eqref{diag4}. In fact its image in $\mathcal{I}(V_{1})$ coincides with 
 \begin{enumerate}
\item $\frac{2}{a_{2}}\cdot \left(T^2- \overline{\left(\frac{a_{2}}{2}+ \frac{r-2}{2\alpha}\right)}T + 1\right)\left[1, X\right]$ \enspace   if   $\tau\leq t-1$ and
\item $\frac{2}{a_{2}} \cdot T\left[1,X\right]$ \enspace  if  $\tau>t-1.$
\end{enumerate}
 Since $\frac{2}{a_{2}}$ is a unit, the operator $T$ is $G$-equivariant and $\mathcal{I}(V_{1})$ is generated by the element $[1,X]$ over $G$, we get that $\sF'_{2}$ is a quotient of $\dfrac{\mathcal{I}(V_{1})}{(T^2-c_{0}T+1)}$ if $\tau \leq t-1$ and $\dfrac{\mathcal{I}(V_{1})}{(T)}$ if $\tau>t-1$.
\end{proof}

\end{document}
